\section{相诊断量}
\label{sec:markers}

\subsection{结构因子与交错关联}
\label{sec:sq}

以 Pauli 算符约定定义静态结构因子
\begin{equation}
S(q)=\frac1L\sum_{i,j}e^{iq(i-j)}\langle Z_iZ_j\rangle.
\label{eq:structure-factor}
\end{equation}
它是使用 $S_i^z=Z_i/2$ 写出同一表达式时的四倍，逐格点关联形式直接适用于开链。
对于测得的比特串 $z_j=\pm1$，单次采样值为
\begin{equation}
x_q=\frac1L\left|\sum_j e^{iqj}z_j\right|^2.
\label{eq:structure-shot}
\end{equation}
其期望等于式~\eqref{eq:structure-factor}，由此可见结构因子非负，且所有 $q$ 均可从同一批采样估计。

$q=\pi$ 时的固定尺度量等于交错磁化的二阶矩：
\begin{equation}
\hat S(\pi)=\frac{S(\pi)}{L}
=\left\langle\left[\frac1L\sum_j(-1)^jZ_j\right]^2\right\rangle.
\label{eq:staggered-moment}
\end{equation}
即使自旋翻转对称性使一阶矩为零，该量仍能反映交错关联。
交错磁化二阶矩是标准的有限尺寸 AFM 诊断量~\cite{sandvik2010computational}。

图~\ref{fig:D03} 在包含 $\pi$ 的同一 $q$ 网格上，对比三个有限尺寸区域各一代表点的 $\hat S(q)=S(q)/L$。
$p=1$ 真机曲线呈现三个代表点特有的 $q$ 依赖，包括 $q=\pi$ 处的 AFM 峰。
后续关联分析采用其中的 $\pi$ 分量。

\begin{figure}[htbp]
\centering
\includegraphics[width=\columnwidth]{{../../figs/exp02/exp02_D03}}
\caption{三代表点的归一化静态结构因子 $\hat S(q)=S(q)/L$：
(a) 数值基态；(b) 同点真机 $p=1$ 测量，使用 $1024$ shots 并经读出缓解。
$q=\pi$ 处峰值诊断交错关联。
\label{fig:D03}}
\end{figure}
